% 第6回　産学連携キックオフ
\session{6}{2}{11月4日（水）}{産学連携キックオフ}{AIとペルソナ設定}{yes}

\goals{
  \item 企業の課題を、背景、制約、期待されている成果に分けて整理できる
  \item ペルソナとは何か、なぜ設計に必要かを説明できる
  \item AIで複数のペルソナ案を作り、現実味と課題との関連で取捨選択できる
  \item 企業から提供された情報の取り扱いを理解する
}

\begin{caution}{企業の情報の取り扱い}
非公開の資料やデータは、AIに入力しない・SNSや外部に書かない・グループ外に共有しない。判断に迷う場合は担当教員に確認する。\quad □ 確認した
\end{caution}

\work{事前準備}{課題提示で聞きたいこと（最低1つ）}
\writelines{2}

\work[80mm]{ワーク1}{課題を整理する}
\begin{wstab}{colspec={Q[l,wd=30mm,m] X[l,m]},row{1-Z}={ht=20mm}}
  \hc{背景\newline なぜこの課題か} & \\
  \hc{制約\newline 予算・敷地・\newline ブランドなど} & \\
  \hc{期待される成果} & \\
  \hc{わからない点\newline 確認したいこと} & \\
\end{wstab}

\work[50mm]{グループ}{メンバーと役割分担}
\begin{wstabh}{colspec={X[2,l,m] X[2,l,m] X[3,l,m]},row{2-Z}={ht=8mm}}
  氏名 & 役割 & 連絡方法・メモ \\
  & & \\ & & \\ & & \\ & & \\ & & \\
\end{wstabh}

\work[50mm]{ワーク2}{AIでペルソナを作り、2〜3体に絞る}
\begin{promptbox}[ペルソナの案を出す]
自動車ディーラーの店舗設計を提案するために、来店客のペルソナを3種類作ってください。各ペルソナには、年齢、職業、家族構成、居住地の特徴、来店の目的、車選びで重視する点、来店時に感じる不安や面倒、休日の過ごし方を含めてください。互いに大きく異なる3人にしてください。
\end{promptbox}
\begin{promptbox}[ペルソナの現実味を確かめる]
次のペルソナについて、現実の顧客として不自然な点、ステレオタイプになっている点を指摘してください。また、このペルソナが店舗に求めることを、本人の言葉で3つ書いてください。（ペルソナを貼る）
\end{promptbox}
\begin{wstabh}{colspec={X[2,l,m] X[3,l,m] Q[c,wd=14mm,m]},row{2-Z}={ht=10mm}}
  ペルソナ案（一言で） & 不自然な点・ステレオタイプ・課題との関連 & 採用\newline ○× \\
  & & \\ & & \\ & & \\ & & \\ & & \\
\end{wstabh}

\work[100mm]{確定}{ペルソナカード}
\noindent
\begin{minipage}[t]{0.49\linewidth}\personacard{A}\end{minipage}\hfill
\begin{minipage}[t]{0.49\linewidth}\personacard{B}\end{minipage}\par

\begin{nexttime}
  \item グループでペルソナ2〜3体を確定し、1枚にまとめる
  \item 課題に関して調べておきたいこと（市場、地域、競合）を分担して集める
  \item AI活用記録を、グループの作業でも続ける方法を決める
\end{nexttime}
\begin{wstabh}{colspec={Q[l,wd=20mm,m] X[3,l,m] X[1,l,m]},row{2-Z}={ht=9mm}}
  分野 & 調べること & 担当 \\
  市場 & & \\ 地域 & & \\ 競合 & & \\
\end{wstabh}
\aimemo
